\subsection{The period entries without a neighbor operator}
\label{subsec:unserved-period}
\hypertarget{source-unserved-period}{}

We first apply the jet and normalization estimates to entries for which no neighboring lattice is inserted. These orders provide the comparison for the entries modified by Hecke operators later in the chapter.

\begin{established}[Places and diagram sizes for the period estimate]\label{hyp:unserved-level-free}
At stage \(j\) use the valuation \(v=v_w/e_w\) for \(w\mid p\),
and let \(H_z\) be the number of boxes at the real label \(z\) with
positive stage weight \(w_z\). Set \(\kappa=k-1\), \(M=N_m\).
For the reference row and the row assigned to \(w\),
\[
 \begin{gathered}
 N_{j-1},N_j\le M,\quad 2\mid N_{j-1},N_j,\quad X''\le a,\quad
 \kappa w_z\le H_z\le a^2,\\
 0\le H_z-\kappa w_z\le4w_z\quad (z\text{ outside the tops}).
 \end{gathered}
\]
Call the place level free for \(\mu\) if
\(p\notin\{T_u,g_0,g_1\}\cup Q(\mu)\), with \(Q(\mu)\) the underlying
twist primes. The ordinary calculation here is restricted to
\(p\nmid N_{\rm host}g_0g_1\) and \(p\le X''\).
\end{established}
\begin{proof}
In \(N_i=k+i(t+1)\) the integer \(k\) is even and \(t\) is odd.
Also \(X''/a=nX'/b\to0\). The real labels are the \(f\) pairs
\((\pm,t)\) of the weight construction. Write \(e=w_z\) and \(N=N_m\). The head coordinates and rank shifts
are those calculated before the integral-kernel comparison in
Proposition~\ref{prop:level-four-integral-kernel}. Off the tops, if \(u\) ranges over the interior Harish--Chandra entries, then
\[
 u_{\max}=4(k-2)\Delta+4k\Delta+\tfrac12,\qquad
 u_{\min}=-\tfrac12,\qquad |u|\le4k\Delta .
\] Thus each of the \(k\)
head rows differs from \(e\) by at most \(S_c+2N\), where one may take
\(S_c=16k\Delta\). Each of the at most \(m(t+1)\) tail rows lies
between zero and \(2\eta\). With the common error
\[
 U=k(S_c+2N)+2m(t+1)\eta,\qquad |H_z-ke|\le U
          \quad(z\text{ off the tops}) .
\]

Here the signed stage offsets and appended rows use the same weight
coordinates as the explicit degree sum in Section~\ref{subsec:selected-exponent}.
On the selected range \(e\ge w''=b/(nX')\), while
\(U\le x^{O(1)}\exp(x^{\epsilon/4})\). Therefore
\[
 U/e\le x^{O(1)}
       \exp(x^{\epsilon/4}+x^{\epsilon/3}-x^\epsilon)\longrightarrow0 .
\]

This proves \((k-1)e\le H_z\le(k+1)e\) off the tops. At a top, delete
the exceptional first row. The \(k-2\) interior head rows have the same
bound as above and all the remaining rows are nonnegative. The deleted
row has at least \(a_t+e-N\) boxes. Since \(a_t\ge xb\) and \(e\le3b\),
for large \(x\)
\[
 H_z-\operatorname{row}_1\ge(k-2)e-U\ge(k-3)e,\qquad
 \operatorname{row}_1\ge2e .
\]

The weights are \(O(a)\) and the number of rows grows polynomially
in \(x\), so their total is at most \(a^2\) beyond one common threshold.
The assigned tuple avoids its coefficient prime when its twist places
are chosen; this is the prime avoidance in the construction of
Proposition~\ref{prop:translate-completion}. The two type primes are
distinct from the discriminant prime.
\end{proof}

\paragraph{The finite sum on a class pair.}
\begin{proposition}[Evaluation of the finite kernel on a pair of classes]\label{prop:unserved-reader-integral}
At a level-free place of Lemma~\ref{hyp:unserved-level-free}, the value of the finite kernel on this pair of classes is the sum
\eqref{eq:unserved-reader}. For each coefficient \(a\) and translation point \(P\) occurring there,
\[
 v_w(a)\ge0,\qquad \gcd(\operatorname{ord}P,p)=1,
\]
  and
the Taylor coefficients in the prescribed differential frames
can be transported coefficientwise by the embedding at \(w\).
\end{proposition}

\begin{proof}
\proofstep{Use the local levels of the two self-dual companions.}

Fix both classes of the matrix entry. Write \(L,L'\) for their
companion lattices in the two standard Hermitian spaces. They are
self-dual by the localization argument of
Proposition~\ref{prop:incident-companion-lattices}. Let \(\Lambda\) be the
embedded period lattice obtained from the inverse half different
times \(\operatorname{Hom}_{\mathcal O_E}(L,L')\), with the
curve period scale. At the finite test places the two half lattices
\(A,B\) have opposite level exponents.

If \(c\) is the positive
integer defining the level of the finite kernel, its prime divisors
are in \(S=\{T_u,g_0,g_1\}\cup Q(\mu)\), and
\[
 cA+cB\subseteq\Lambda,\qquad c\Lambda\subseteq A\cap B.
\]
These follow by comparing each of the exponents with the
valuation of \(c\) locally, including the zero exponents off the
level. Write \(S_X,S_A,S_B\) for representatives, respectively, of
\[
 A/(A\cap\Lambda),\quad
 (\Lambda\cap A)/(\Lambda\cap A\cap B),\quad
 B/(B\cap\Lambda).
\]
In particular \(h=|S_A|\) is prime to \(p\). Indeed the middle
quotient is finite and killed by \(c\).

\proofstep{Compute the finite Fourier expansion.}

Put \(E_H(x,y)=\Im H(x,y)\) for the alternating period form,
\(\alpha(x)=\exp(\pi i\sqrt d\, H(x,x)/2)\), and let
\(\Xi\) be the actual finite test extended by zero off the
embedded rational Hom space. For the normalized theta
\(\Theta_{L,L'}\), the finite part of the entry is
\begin{equation}\label{eq:unserved-reader}
 \sum_{x\in S_X,\ y\in S_B}
   \alpha(x)\alpha(y)e^{-\pi i E_H(x,y)}
   \left(\frac1h\sum_{s\in S_A}
      \Xi(s-x)e^{2\pi i E_H(s,y)}\right) U_{x+y}\Theta_{L,L'} .
\end{equation}
The Fourier sum at \(y\) uses the inverse commutator \(\exp(-2\pi i E_H(s,y))\), with the counting
measure divided by \(h\). The kernel reading uses volume one.
Applying \(U_x\) to the resulting \(U_y\) gives the displayed
Weyl factor. Since \(E_H(\Lambda,\Lambda)\subset\mathbb Z\)
and \(\sqrt d\,H(v,v)/2\in\mathbb Z\) on \(\Lambda\), all three
phase factors and the Fourier character are roots of unity.
The translates have order dividing \(c\) modulo \(\Lambda\).

\proofstep{Check the additive and determinant character factors.}

The word support is
a characteristic function. The type tests and twist tests on
a residue matrix multiply finite additive characters and values
of finite multiplicative characters, with zero on the excluded
residue classes. Each nonzero value is a root of unity.
Thus, for a nonzero finite character factor \(\xi\) of order \(n\),
\[
  n\,v_w(\xi)=v_w(\xi^n)=v_w(1)=0 .
\]
 In a
moved test the extra factors are integer powers of square roots
of the residue cardinality and local determinant characters.
At a type place the two movers preserve the submitted lattices,
so their determinants are local units and the residue characters
suffice.

At a twist place the uniformizer value and its inverse
are integral away from that place's rational prime and the type
levels. Indeed raise the ideal of the place to a principal power
with generator congruent to one at the character conductor.
The defining algebraic infinity type expresses that power of
the value as a product of conjugates of the generator to
integer powers, up to a root of unity. If \(\mathfrak P^A=(\alpha)\) with \(\alpha\equiv1\)
at the conductor, the value \(u\) satisfies

\[
  u^A=\zeta\prod_\tau\tau(\alpha)^{a_\tau}
\]
in the algebraic normalization. Thus
\[
       A v_w(u)=\sum_\tau a_\tau v_w(\tau(\alpha))=0
\]
for \(w\) outside that rational prime and the type levels. The off-type matching of the
column character with this ideal character gives the assertion
for the actual local value. The types are disjoint from the
auxiliary twist places. Finally the pinned terminal row either
uses an indicator or its finite Fourier reading. Its nonzero
values are integral additive character values, with a possible
factor \(q^{-Mn}\) at a place of cardinality \(q\), rank \(n\)
and depth \(M\). At the coefficient place under consideration, \(q\) is a unit. For any nonzero additive-character value \(\xi\) this gives
\[
  v_w(q^{-Mn}\xi)=-Mn\,v_w(q)+v_w(\xi)=0.
\]
 The
nonterminal classes use the matrix test, by the terminal pinning
of the supplementary rows. This proves integrality of every
scalar multiplying a translate in \eqref{eq:unserved-reader}.
It applies also when an entire entry vanishes.
\label{cor:unserved-kernel-charts-1}

All Taylor identities here are identities of coefficients.
\proofstep{Transport the Taylor coefficients in their fixed normalization.}

For clarity, algebraicity of the translated derivatives can
also be tested at one auxiliary good prime. Choose it away from
the denominators of the period comparison, the translate and a
rational basis of the Hom directions, away from the discriminant,
and split completely in the host. The fourfold expansion and the
baseline derivative estimate of Section~\ref{subsec:baseline-jets}
apply for every isomorphism of \(\mathbb C\) with \(\mathbb C_p\)
over a fixed embedding of \(\overline{\mathbb Q}\) at this auxiliary
prime. For each fixed derivative \(c\) in those directions the resulting bound \(M_c\) is uniform in the chosen isomorphism:
\[
  |\iota(c)|_p\le M_c\qquad\bigl(\iota:\mathbb C\xrightarrow{\sim}\mathbb C_p\text{ over the fixed algebraic embedding}\bigr).
\]
 In particular, the argument matrices are
integral after a fixed denominator clear for the rational
directions. The embedded matrices of a rational basis over the
quadratic field span the complexified Hom at the CM type, by
separability of \(E/K\).

For a fixed total order \(d\), the induced action on derivatives
is the degree-\(d\) tensor of the inverse direction matrix. Its
entries are algebraic and fixed. Thus every labelled coefficient
is a fixed finite algebraic linear combination of the bounded
directional derivatives. The triangle inequality gives the same
uniformity in the field isomorphisms.
Suppose one coefficient \(c\) were transcendental, and fix one
of the comparison isomorphisms \(\iota\colon\mathbb C\to\mathbb C_p\).
For each \(n\ge1\), the map 
\[
  c\mapsto p^{-n}c
\]
 extends by
transcendence bases to an automorphism \(\alpha_n\) over
\(\overline{\mathbb Q}\). The same directional bound would then
apply to all the isomorphisms \(\iota\circ\alpha_n\), whereas
\[
 |(\iota\circ\alpha_n)(c)|_p=p^n|\iota(c)|_p\longrightarrow\infty .
\] Division by the multi-index factorial preserves
algebraicity. The finite sums, linear substitutions and the
map on derivatives can consequently be transported by the
specified coefficient embedding at the original \(w\).

\end{proof}

\paragraph{Removing the differential generator in the actual frames.}

A differential generator can change the valuation in a scalar chart. In the actual integral frames the following comparison removes that factor.

\begin{proposition}[Removing the generator in integral frames]\label{prop:unserved-normalized-frame}
At the enhanced places in this comparison, including \(p=d\),
one can choose the class frames so that the substitution of \eqref{eq:unserved-normalized-frame} reads
\[
 Z_\tau\longmapsto \rho_\tau(P_{i+1})Z_\tau\rho_\tau(P_i)^{-1}.
\]
 The axis
generators and the resulting scalar factors in the coefficient
tests are units in the chosen coefficient valuation.
\end{proposition}

\begin{proof}
\proofstep{Choose simultaneous rational unitary frames.}

For each stage \(i\) and each of its classes choose a rational
unitary matrix \(P_i\) whose columns span its lattice at all
places above \(p\). Such a matrix exists simultaneously there.
The selected host extension \(E/F\) is unramified at finite places.
At a split place the two summands of a self-dual lattice are
mutually dual.
In paired local coordinates the Gram block between dual bases \(e_i,e'_j\) is
\[
  \big(h(e_i,e'_j)\big)_{ij}=I .
\]

At an unramified nonsplit place, orthogonal
splitting and surjectivity of the unramified unit norm identify
two unimodular Hermitian lattices. The resulting local
transporters are open. Multiplication by nearby scalar norm-one
elements avoids the eigenvalue \(-1\) at each place. Approximate
their inverse Cayley transforms in the rational vector space of
skew-Hermitian matrices by weak approximation over \(F\), and
apply the Cayley transform back. Thus both the columns and the
inverse coordinates clear into the required lattices by
integers prime to \(p\). The same matrix \(P\) works at every label over the place. Its inverse adjoint property is
\[
  P=(P^*)^{-1},\qquad P^*P=I.
\]

\proofstep{Separate the residual axes.}

For a CM type embedding \(\rho_\tau\) at \(w\) write \(d_\tau\)
for the image of the generator of its axis in the module
spanned by \((\rho_z a)_z\), \(a\) in the inverse half different.
At these enhanced places that ideal is the unit ideal locally.
The entries are integral. For two distinct type labels \(\tau,\tau'\), one can consequently separate their residue embeddings by an integer:
\[
  a\in{\cal O}_E,\qquad \overline{\rho_\tau(a)}\ne\overline{\rho_{\tau'}(a)}.
\]
 Indeed the labels
all extend the same embedding of \(K\), and the relative inertia
in \(E/K\) is trivial at these primes.

Distinct primes are
separated by the Chinese remainder theorem and distinct
embeddings of the same unramified local field are separated
on the residue field. This also applies at \(p=d\), where
\(K/\mathbb Q\) itself is ramified. Acting on the embedded
all-ones vector with the product, over \(z\ne\tau\), of
\(a_z-\rho_z(a_z)\), where
\(\rho_\tau(a_z)-\rho_z(a_z)\) is a unit, isolates the
\(\tau\)-axis with unit coordinate. The module is stable under
these actions. Its axis generator is therefore a unit, and
membership of the generator supplies a finite expression
\[
 d_\tau {\bf1}_\tau=\sum_l r_l(\rho_z a_l)_z,
 \qquad |r_l|_w\le1,\quad a_l\in\mathfrak b^{-1},
\]
where \(\mathfrak b\) denotes the half different.
This last expression keeps its global period directions.

\proofstep{Absorb the unit without changing the norm tests.}

Set \(a_\tau=d_\tau^{-1}\) and use the frame
\(a_\tau^i\rho_\tau(P_i)\) at stage \(i\). In the incident
Hom variable the factor
\(d_\tau a_\tau^{i+1}/a_\tau^i\) is exactly one. The substitution
on the series in \eqref{eq:unserved-reader} is now
\begin{equation}\label{eq:unserved-normalized-frame}
 Z_\tau\longmapsto
    \rho_\tau(P_{i+1}) Z_\tau\rho_\tau(P_i)^{-1},
\end{equation}
with the row of \(Z_\tau\) indexed by the target. These are
admissible choices for the true columns.

The product of the
column minors of a diagram \(\nu_\tau\) gains
\((a_\tau^i)^{|\nu_\tau|}\) on input frames. A determinant twist
of exponent \(s\) gains \((a_\tau^{iN_i})^s\), including signed
exponents. The target formulas use \(i+1\). Every such factor
is a unit, and homogeneity gives the same power on every term
of the corresponding polynomial reading. Thus the original
coefficient norm tests survive this change of frames, while the
Taylor map itself has the precise coefficient in
\eqref{eq:unserved-normalized-frame}. \label{cor:unserved-kernel-charts-2}

\end{proof}

\paragraph{One curve model before the inert orbit estimates.}
\begin{proposition}[One inert curve model for all translates]\label{prop:unserved-common-inert-model}
At the selected inert places \(p\ge5\) with \(p\ne d\), the two classes and all their prime-to-\(p\) translates use one good curve law, tangent unit \(c_w\) and quadratic correction \(q\). In good curve coordinates with logarithm \(l\), the parameter change for all these charts is
\[
 t(X)=l^{-1}(c_w L(X)),\qquad c_w\in{\cal O}_w^\times.
\]

In the even and signal cases the resulting normal form is
\textup{\eqref{eq:tag-i4}}. The odd case uses the doubled orbit and
the zero-step case uses the repeated-variable one-letter
phase. The phase and its coefficient estimates are uniform
in the labels, although the integral series may depend
on the classes and on every variable.
\end{proposition}

\begin{proof}
\proofstep{Fix the model and the common unit.}

The four cases treated here share the following hypotheses on the residue characteristic
\[
  p\ge5,\qquad p\ne d,\qquad \left(\frac{-d}{p}\right)=-1.
\]
Apply Proposition~\ref{prop:good-inert-cm-model} to the discriminant-normalized
period lattice. Changing the integral ideal representative of the
curve prime to \(p\) changes the discriminant root by the inverse
homothety and leaves its scaled lattice fixed. The good-model
homothety is a unit since its twelfth power compares two unit
discriminants. Write \(l\) for the good curve logarithm and \(L\) for the comparison logarithm. The integral parameter change \(t\) fixes the tangent unit through
\[
  l(t(X))=c_wL(X),\qquad t'(0)=c_w\in{\cal O}_w^\times.
\]
Use this one choice of parameter, Frobenius and curve phase \(q\) throughout the place. These precede the choice of the two classes, their ranks,
their diagrams, and the translations in the reader.

\proofstep{Choose the underlying Hom frame before the translates.}

Here is a rational orthonormal chart for the Hom of the actual
pair. For the standard lattices one matrix coordinate carries
the trace Hermitian form on the inverse half different over \(K\).
It is unimodular at \(p\) by duality of the trace lattice
(the half different times its conjugate is the relative
different).
Split off a vector of unit norm from its reduction.
For a lift \(u\) with unit norm, let \(\ell\) be its Hermitian pairing in the linear variable. The integral projection and its kernel are
\[
  x\longmapsto \ell(x)\ell(u)^{-1}u,\qquad \ell(u)\in{\cal O}^{\times},\qquad \ker\ell=u^\perp.
\]
 Such a
vector exists in every nonzero nonsingular Hermitian space over
the quadratic residue field, by polarization. Lift it to the
localized rational lattice.

Orthogonal projection onto its
complement divides only by that unit, so iteration gives an
orthogonal basis with positive rational norms \(n_j/e_j\) whose
numerators and denominators are prime to \(p\). Positivity and
rationality follow from definiteness and the fixed field of
conjugation on \(K\). In four copies of each line choose four
integer squares summing to \(n_je_j\). If \(Q_j\) is their quaternion matrix, then
\[
 Q_j^tQ_j=n_je_j I_4,\qquad
 (Q_j/n_j)^t\bigl((n_j/e_j)I_4\bigr)(Q_j/n_j)=I_4 .
\]
Both \(Q_j/n_j\) and its inverse are integral at \(p\). Use these four copies in each
\(2\)-by-\(2\) block of matrix entries, possible since both
ranks are even. Transport the resulting basis by

\[
  X\mapsto P_{i+1} X P_i^{-1}.
\]
 This is a rational isometry
and spans the Hom lattice over the local quadratic integers.
The basis is chosen before any translate.

\proofstep{Descend the corrected integral section from the fourth tensor.}

On four copies the normalized theta expansion of
Section~\ref{subsec:theta-normalization} has integral coefficients
in the locally matched frame. Its arguments are integral
unitary transforms of the orthonormal directions just
constructed. The one-variable translated section multiplied by
\(\exp(q)\) is integral in group coordinates. For the sum of
phases a unitary change preserves its class modulo logarithms
of principal integral units by the polarized CM norm relation
\eqref{eq:cm-norm-relation}.

Apply the change and extract four
equal coefficient indices on disjoint copies. This proves, in the
orthonormal basis \(b_j\), that the product
\[
  \exp\!\left(\sum_j q(X_j)\right)
  \Theta\!\left(\left(c_w\sum_j\rho_\tau(b_j)L(X_j)
                        \right)_\tau\right)
\]
is integral, for each of the normalized translates \(\Theta\)
in the finite reader. It holds also for an integral linear
combination since the multiplier and arguments are common. A pair
of lattices agreeing only locally at \(p\) gives the same conclusion.
Two integers prime to \(p\) clear the source and target lattice
inclusions in both directions, preserving these local arguments
and prime-to-\(p\) torsion.

\proofstep{Read the phases in the trace-dual cycle variables.}

In an even orbit expand the two parities in a basis and its
trace dual over the quadratic local field. For one pair of parity vectors \(e,o\) in a matrix coordinate, the Gram block is
\[
  \begin{pmatrix}h(e,e)&h(e,o)\\ h(o,e)&h(o,o)\end{pmatrix}=\begin{pmatrix}0&1\\1&0\end{pmatrix}.
\]
 This follows by multiplying the unitary
source and target changes in
\eqref{eq:unserved-normalized-frame} by their inverse adjoints
on the conjugate coordinates. Thus the change from the
orthonormal variables gives the hyperbolic substitution of
Section~\ref{subsec:curve-exponents}, with the trace-dual sums
there called \(E\) and \(O\). The root-of-unity average over
\(p^{d_p}-1\) elements is integral and replaces the summed
polarization by \(B_{d_p}\) modulo logarithms of integral
units.
The averaging denominator is a unit, since
\[
  p^{d_p}-1\equiv-1\pmod p,\qquad v_p\big((p^{d_p}-1)^{-1}\big)=0.
\]

Returning to additive coordinates gives \eqref{eq:tag-i4} of
Section~\ref{subsec:incident-sections}, with exponent
\(-B_{d_p}\) on each entry and argument scale given by the
common tangent unit. This works for any choice of the orbit
representatives and includes the signal orbits. The same
orthonormal identity at odd order is used below with the
trace-dual variables on the doubled orbit. At order one the
labelled basis is already unitary over the quadratic integers.
The norm relation transports the one-variable gauge, and its
polarization at the repeated variable gives the one-letter
phase with the factor \(1/2\). Thus in each case the integral
factor can depend on the classes and on every variable, while
the phase and its coefficient estimates are uniform in both
classes and all entries of their diagrams. \label{cor:unserved-kernel-charts-3}
\end{proof}


\subsubsection{Even cycles and signal primes}

On an even cycle the cyclic jet estimate gives a contribution from its paired positions. For a signal prime the same accounting must count each selected cycle only once.

\begin{proposition}[Entries on the even cycles]\label{prop:unserved-even-entries}
In the ordinary cutoff range of Lemma~\ref{hyp:unserved-level-free} with \(p\ge5\), suppose the
inert step has positive even order or the \(p\)-cycles are signal
cycles. In the latter case the cycles are counted once, with \(A\)
empty. For the full entry \(b\), using the even-cycle allowances below,
\[
 v_p(b)\ge\sum_r g_r+F_{\rm in}-E_{\rm I}.
\]
The estimate includes a signal cycle of length two.
\end{proposition}

\begin{proof}
\proofstep{The zero-index estimate on the first cycles.}

A positive even order of the inert Frobenius step gives one
set \(A\) of cycles on the first embedding and one set
\(BC\) on the second. At a signal prime Frobenius changes the
real embedding and the cycles are counted once, with
representatives on the first embedding. Then \(A\) is empty.
In both cases the cycle length \(d_p\) is even and at least
two, and \(p\ge5\). On a cycle in \(A\) the estimate with
index zero in \eqref{eq:tag-j1} implies the gain \(\max(g_0(r),0)\) at its
representative \(r\). In detail, the exclusion at a top
omits each least forward odd distance to a top, hence
the remaining terms in \(g_0(r)\) occur in the first sum of
\eqref{eq:tag-j1}. Its numerator \(\alpha_s\) is at most the diagram's
number of height-one columns and its recipient width
bound is at least the actual remaining width. With

\[
  \Lambda=\log_p(2a)\ge1,\qquad 1+\log_p(2a)\le 2\Lambda,\qquad 1+\log_p(a^2)\le3\Lambda,
\]
 the digit counts for a
height-one count at most \(2a\) and a remaining size at
most \(a^2\) are at most \(2\Lambda\) and \(3\Lambda\).
The error in \eqref{eq:tag-j1} is thus at most
\[
   5d_p M^2\Lambda+4\,\#\{\text{tops on }r\}\Lambda^2 .
\]
This is covered by the digit terms in \(g_0(r)\).
The nonnegative alternative in \eqref{eq:tag-j1} gives the maximum
with zero. This argument uses the zero-index coefficient
of the same homogeneous jet and is independent of the
choice of the neighbor summands on any other cycles.

\proofstep{Cyclic flow and the top capacity.}

For \(BC\) use the cyclic bound of
Proposition~\ref{prop:cyclic-jet} with the stage weights.
The capacity comparison \eqref{eq:flow-comparison} gives
the cyclic flow term 
\[
  F_r={\cal F}_p((\kappa w_z)_{z\in r})
\]
and the top contribution. At the second embedding there
are no tops and the top contribution is zero. On a signal
cycle the top contribution pays \(\max(g_0(r),0)\) after
the top error. Denote the resulting error by
\(\operatorname{err}_r\). Put \(n_r=\#\{\text{tops on }r\}\)
and \(W_*=b+\eta_{\rm sz}+1+3S_{\rm sz}\), the uniform
upper bound for the remaining width at a top. The error is
\[
 \begin{split}
 \operatorname{err}_r={}&\gamma n_r W_*+
  \frac{c_6}{p}\sum_{z\in r}w_z+
  d_p N_{j-1}N_j\bigl(1+\log_p(4d_p M a)\bigr)\\
 &+K_1d_p M^2\log_p(2a)+K_2 n_r\log_p^2(2a).
 \end{split}
\]
The constants satisfy \(c_6\ge12\), \(K_1\ge5\),
\(K_2\ge4\). The first two terms dominate
the top width and the loss \(11\sum w_z/p\) in
\eqref{eq:flow-comparison}.

The third and fourth terms cover
\(d_p N_{j-1}N_j(1+\log_p(a^2))\), since \(4d_pM\ge1\)
and \(K_1 M^2\log_p(2a)\ge N_{j-1}N_j\log_p a\).
We obtain the bound
\[
  \begin{cases}
    g_r, &r\in A,\\
    g_r+F_r-\operatorname{err}_r, &r\in BC,
  \end{cases}
\]
where \(g_r=\max(g_0(r),0)\) on first-embedding representatives
and is zero on second-embedding representatives. On the signal
cycles each odd step goes to the second embedding, which has
no tops. The distance in \(T_s\) is therefore infinite.
More explicitly, for any first-embedding top \(s\) on a signal cycle one has
\[
  \{j>0:j\ {\rm odd},\ s+j\text{ is a top}\}=\varnothing .
\]
The recipient width bound and the lower bound on the top's
height-one count show that \(T_s\) dominates the positive-part
sum in \(g_0\) term by term. That sum before the digit correction
is nonnegative. This pays \(g_r\) before the other errors in
\(\operatorname{err}_r\).

\proofstep{Add the bounds over the actual orbits.}

The logarithm and its inverse preserve disjoint sets of orbit
variables. The integral factor of the full theta series can
depend on all orbits. Expand it in monomials before inverting
the logarithm. Each monomial gives a product of admissible
single-orbit factors, and only finitely many meet a specified
jet. The derivative polynomial is a product over the orbits
on such a term.
For a monomial product \(\prod_r H_r(Z_r)\) in disjoint orbit variables the jet factors as
\[
  \partial^\alpha\!\left(\prod_r H_r\right)(0)=\prod_r\partial^{\alpha_r}H_r(0),\qquad\alpha=(\alpha_r)_r.
\]
 The valuation bounds therefore add. This
also explains the restrictions to a single orbit: setting the
other variables equal to zero commutes with their logarithms
and deletes their phase, whose constant is zero. The common
scaling of logarithmic coordinates is by a unit.

Each orbit term \(F_r\) is invariant under rotation, so
\(F_{\rm in}=\sum_r F_r\) is also the sum over all labels of
the same orbit terms divided by \(d_p\). The error
\(E_{\rm I}\) includes \(\sum_{BC}\operatorname{err}_r+
\max(0,\sum_A F_r)\). The summed entry estimate is therefore
at least \(\sum_r g_r+F_{\rm in}-E_{\rm I}\). In particular
the real-embedding count here is also valid for a signal
orbit of length two.

\end{proof}

\subsubsection{Odd cycle lengths}

An odd cycle has no alternating pairing of all its positions. Its contribution is estimated from the forced exponent and the ordinary error for this orbit.

\begin{proposition}[Entries on the odd cycles]\label{prop:unserved-odd-entries}
For \(p\ge5\) in the ordinary cutoff range of Lemma~\ref{hyp:unserved-level-free}, let \(b\) be the entry at a positive odd inert step of order \(d_p\ge3\). Then
\[
 v_p(b)\ge F_{\rm in}-E_{\rm I}.
\]
 There is no additional \(g_0\) gain
in this case. The phase comparison uses the doubled even cycle,
and the flow error includes \(\kappa\sum_{BC}w_z/p^3\).
\end{proposition}

\begin{proof}
\proofstep{Realize the phase on the doubled cycle.}

Here \(d_p\ge3\) is odd and the two sets of cycles again lie
over separate real embeddings. In the standard lattice
bases the local chart has an integral factor in the group
coordinates and an exponential of the sum of orbit phases
in additive coordinates. We explain the chart comparison
because the actual orbit has odd order and the parity in the
phase calculation is on a different set of letters.

Write \(\sigma\) for Frobenius on the unramified local
coefficients. Take an integral basis \(\omega_\rho\) of the
unramified extension of degree \(d_p\) over the unramified
quadratic field, orthonormal for the trace Hermitian form
whose conjugation is \(\sigma^{d_p}\).
The two parities belong to the doubled indexing set: reducing positions to the actual odd cycle identifies
\[
  s,\ s+d_p\ \in\mathbb Z/(2d_p)\quad\longmapsto\quad s\bmod d_p,\qquad (-1)^{s+d_p}=-(-1)^s.
\]
 This form is
unimodular by the unramified trace pairing. Reduction,
orthogonal splitting and surjectivity of the norm on
quadratic units give the basis, as in
Section~\ref{subsec:theta-normalization}. Its trace-dual basis
is \(\omega_\rho^*=\sigma^{d_p}\omega_\rho\).
In these coordinates trace orthogonality is
\[
  \operatorname{Tr}\big(\omega_\rho^*\omega_{\rho'}\big)=\delta_{\rho\rho'}.
\]

On a cycle with representative \(r\) and fixed matrix entry,
the trace-coordinate linear variables are
\[
 L_{r,\rho}(Z)=\sum_{\substack{0\le s<2d_p\\s\ {\rm even}}}
              \sigma^s(\omega_\rho^*) Z_{r+s}.
\]
Here \(r+s\) refers to the actual orbit of size \(d_p\).
The even labels give every position once. Integral standard
bases for the two incident lattices give integral Frobenius
coordinates \(g_s\) with 
\[
  \sigma^2 g_s=g_{s+2}.
\]
 Their
conjugate components give the paired evaluations. The
coefficients 
\[
  \sum_{s\ {\rm even}}\sigma^s(\omega_\rho)g_s
\]
are fixed by \(\sigma^2\), and the trace-dual identity gives
the inverse expansion of \(g_s\) in the displayed linear
variables. In the orthonormal kernel frame the source and
target evaluations are in perfect duality, entry by entry.
Expanding the Hermitian Gram matrix of this change of
variables and using trace duality thus gives identity:
it is zero between distinct orbits or entries, and is
\(\delta_{\rho\rho'}\) in each such block. It follows that
this linear change preserves the integral exponential
comparison in the orthogonal chart.

In group coordinates the backward steps of length \(2e\)
in the generalized logarithm give the same \(L_{r,\rho}\)
after applying the one-dimensional logarithm to the
integral formal sum with the indicated coefficients.
Write \(\mathcal G_{r,\rho}\) for that logarithm in the group coordinates. For the even representative \(s\), its coefficient is
\[
  [Z_{r+s}^{p^{2e}}]\,\mathcal G_{r,\rho}=p^{-e}\sigma^{s+2e}(\omega_\rho^*).
\]
 This is the Frobenius-twisted
logarithm of the formal sum by \eqref{eq:integral-log}.

The even and odd
linear forms in the doubled cycle agree on the two
identified copies when the odd form uses \(\omega_\rho\):
translation by \(d_p\) swaps their parity and applies
\(\sigma^{d_p}\) to the coefficient. The norm relation
therefore takes the sum of the diagonal phases to one
half of the averaged two-slot phase, modulo a phase
whose exponential is integral in group coordinates.
Here division by two comes from polarization, and
preserves this congruence because \(p\) is odd. In the
doubled cycle the averaged phase has the hypotheses of
Lemma~\ref{lem:two-differences}.

\proofstep{Pay the exceptional weights by the odd-step inequalities.}

For a nonnegative family \(b_z\) satisfying
\[
              p^q b_z+b_{z+q}\ge1\qquad(q>0\text{ odd})
\]
the phase at weights \(\rho+b_z\) has weight at least
\(\theta\). The exceptional monomials in the doubled
cycle project to two entries related by an odd step,
possibly to the same label. The displayed inequality pays
their missing unit of weight.

If both exponents are powers
of \(p\) with a common positive power, removing that
common power reduces the constraint to one of the displayed
inequalities. The odd step is measured before reducing labels
modulo \(d_p\). The inverse logarithm has weight at least
\(\rho\) in every coordinate, by the same induction as for
the even logarithm. It follows that the full rescaled
series has weight at least zero. Thus, including the
factorials, the valuation of a derivative of diagram size
\(H_z\) at every label is bounded below by
\[
          \sum_z H_z(\gamma-b_z)-
                    f N_{j-1}N_j(1+\log_p(a^2)).
\]
One can use this monomial estimate before applying the
integral coefficients of the framed minors.

\proofstep{Choose the dual weights on the two embeddings.}

On first-embedding cycles use \(b_z=1/(p+1)\). This is
admissible and gives a nonnegative contribution to the first
sum. On the second embedding choose an optimum dual for
the finite program with capacities \(\kappa w_z\) and
positive odd steps strictly less than \(d_p\).

Choose it in
\([0,1]\) and raise each coordinate to at least \(p^{-3}\).
This adds all omitted constraints, since any omitted odd
step is at least three. The cost on the original capacities
is at most \(\kappa\sum w_z/p^3\).
Writing \(b'_z=\max(b_z,p^{-3})\) for this modification, one obtains
\[
  \sum_z\kappa w_z(b'_z-b_z)\le p^{-3}\kappa\sum_z w_z .
\]
 A vertex at which the
original dual is at least \(2/p\) has saturated capacity
and no outgoing flow, by complementary slackness.
The sum of its weights over these vertices is at most
\(\sum w_z/p\), by the total source capacity of the flow.

The increase from \(\kappa w_z\) to \(H_z\) can therefore
reduce the lower bound by at most \(8\sum w_z/p\).
Indeed use \(b_z\le1\) on those vertices and
\(\gamma-b_z\ge-1/p\) elsewhere. The digit term is at most
\(3fM^2\log_p a\) since \(p\le a\). The resulting bound is
\[
 \sum_{r\in BC}F_r-
     \frac{\kappa}{p^3}\sum_{BC}w_z-
     \frac8p\sum_{BC}w_z-3fM^2\log_p a .
\]
Here \(E_{\rm I}\) contains the additional term
\(\kappa\sum_{BC}w_z/p^3\), as well as the other
allowances for these losses. There are no tops in \(BC\).
Subtracting \(\max(0,\sum_A F_r)\) covers the entire
first-embedding flow term. Thus this bound is at least
\(F_{\rm in}-E_{\rm I}\). The operator-free gain \(g_0\)
is zero in the assembled estimate at these primes.

\end{proof}

\subsubsection{The one-letter and ramified cases}

A zero inert step leaves one letter at each real label. This reduces its contribution to the one-letter derivative estimate and leaves the ramified primes to be treated separately.

\begin{proposition}[Entries at a zero inert step]\label{prop:unserved-zero-entries}
At a zero inert step with \(p\ge5\) and the ordinary prime range
and diagram sizes of Lemma~\ref{hyp:unserved-level-free}, there
is one letter on each real label. By the jet estimate of
Lemma~\ref{lem:one-letter-derivatives}, the entry \(b\) satisfies
\[
 v_p(b)\ge F_{\rm in}-E_{\rm C}.
\]

\end{proposition}

\begin{proof}
A zero inert step gives one letter at each real place.
Its local phase is the two-slot phase evaluated on
\((Z,\eta_z Z)\) and divided by two, with \(\eta_z\)
a unit. The one-letter calculation
\eqref{eq:one-letter-jet} applies separately to each label.
Expanding the common integral factor in group coordinates
as for the even orbits gives the estimate
\[
          \rho\sum_z H_z-f N_{j-1}N_j(1+\log_p(a^2)).
\]
In this case \(F_{\rm in}=\rho\kappa\sum_z w_z\).
The digit allowance in \(E_{\rm C}\) is
\(3 fM^2\log_p a\), and its remaining
\(\rho\sum_z w_z\) is nonnegative. This proves the
needed comparison with \(F_{\rm in}-E_{\rm C}\).

\end{proof}

\begin{proposition}[Entries at the discriminant prime]\label{prop:unserved-discriminant-entries}
At the discriminant prime \(p=d\), use the paired embeddings at
each real label, the level-free test and the normalized frame
\eqref{eq:unserved-normalized-frame}. Write \(b\) for the entry of the word \(u\). The counts below and the common budget give
\[
 v_p(b)\ge-\tfrac12\sum_t c_t(u)\ge F_{\rm D}-E_{\rm D}.
\]
 This comparison has no \(X''\) cutoff.
\end{proposition}

\begin{proof}
\proofstep{Apply the paired jet bound to the diagram.}

At the discriminant, let \(A_t^\pm\) count columns and
let \(a_t^-=A_t^--(A_t^-)'\) count columns of height one
on the first embedding. For the paired quadratic form of Section~\ref{subsec:cycle-jets}, the bound at the ramified prime is \[ -\frac12\sum_t c_t(u),\qquad c_t(u)=H_{(-,t)}-\mathbf1_{t\ {\rm top}}a_t^- +\mathbf1_{t\ {\rm top}}A_t^+ . \] The \(t\)-pairs are the two embeddings of the same real label, including at a top. On a generic
minus-frame change the plus directions change by its
inverse transpose. On the first row use the plus
column count, and on the other rows use the minus
degree, as in the ramified calculation following
Proposition~\ref{prop:cyclic-jet}. The quantity subtracted
from the first-row count is only the number of
height-one columns. The budget defines

\[
  E_{\rm D}=F_{\rm D}+\tfrac12\max_u\sum_t c_t(u).
\]
It follows that each word's entry has valuation at
least \(F_{\rm D}-E_{\rm D}\), with no use of \(X''\).

\proofstep{Transport the curve phase in the ramified formal coordinates.}

Here is the source of the paired shape used at the
discriminant. In the good curve model write \(Z=l(t)\) for
the formal logarithm and \(\phi(Z)=\mathrm{ph}(l^{-1}(Z))\)
for the one-variable phase in additive coordinates. The
formal group law and \(l'(t)\), with \(l'(0)=1\), are
integral. To see the curve congruences at this particular prime,
use the discriminant-one model as in
Proposition~\ref{prop:level-free-theta-jets}; \(d\equiv7\pmod8\) and

\[
  p=d\ge7,
\]
 so a good short equation is available over a finite
extension. The sigma addition calculation of
Section~\ref{subsec:curve-exponents} gives the parallelogram
identity modulo logarithms of principal integral units in the
curve parameter. Its proof cancels the entire \(s_2\) quadratic
and uses only good reduction. Polarization by the integral formal
half therefore gives biadditivity modulo these logarithms. For
\(a=1+\sqrt{-d}\), \(n=a\bar a=1+d\) is a unit. The isogeny
identity \eqref{eq:cm-isogeny} gives

\[
  q_E([a]t)\sim n q_E(t):
\]
 each nonzero kernel point has
prime-to-\(p\) order and integral affine coordinates, and the
leading factor \([a]t/(at)\) is a principal integral unit.
Polarize this identity, substitute \([a]^{-1}\) in the second
argument and use \(n[a]^{-1}=[\bar a]\). Subtract the adjoint
identity for the scalar one by biadditivity to obtain the
adjoint identity for \(\sqrt{-d}\). In particular this computation
does not require the inert unramified formal module at the
discriminant.

Let \(B(Z,W)=\phi(Z+W)-\phi(Z)-\phi(W)\), put
\(Z_i=l(t_i)\), and write
\[
  C=B(Z_1+Z_2,Z_3)-B(Z_1,Z_3)-B(Z_2,Z_3).
\]
The first-level congruence gives \(U=\exp(C)\) an integral
unit. Therefore
\[
 \partial_{t_1} C=U^{-1}\partial_{t_1}U
    \quad\text{in the integral formal series ring}.
\]
The other two derivatives of this series are integral.
In the expansion of \(C\), every summand other than
\(\phi(Z_1+Z_2+Z_3)\) omits at least one variable, so
\[
 \left.\partial_{t_1}\partial_{t_2}\partial_{t_3}C
             \right|_{t_2=t_3=0}
     =l'(t_1)\phi'''(l(t_1)).
\]
Here \(l'(0)=1\), and \(l'\) is an integral unit.
It follows that \(\phi'''(l(t))\) is integral.
Applying \(l'(t)^{-1}d/dt\) repeatedly and evaluating at zero
shows integrality of \(\phi^{(j)}(0)\) for every \(j\ge3\).
The congruence at multiplication by \(-1\) gives
\(2\phi'(0)\) integral.

Apply the same logarithmic-derivative argument in two
variables to the adjoint congruence for
\(B(\alpha Z,W)-B(Z,\bar\alpha W)\) with \(\alpha=\sqrt{-d}\)
and \(p=d\). Its mixed derivative at the origin is
\[
  (\alpha-\bar\alpha)\phi''(0)=2\sqrt{-d}\,S ,
              \qquad S=\phi''(0).
\]
This is integral, and 
\[
  v(2\sqrt{-d})=1/2
\]
 in the normalized
valuation. Hence 
\[
  v(S)\ge-1/2.
\]

\proofstep{Compute the quadratic phase in the orthogonal model.}

In the normalized orthogonal coordinates separate
\(\tfrac12 SZ^2\) from \(\phi(Z)\). The remaining series
has integral derivatives. So does the inverse logarithm.
Write \(l\) for the group logarithm and \(g=l^{-1}\). Its invariant differential has integral coefficient series with constant one, and differentiation gives
\[
  l'(g(Z))\,g'(Z)=1,\qquad g'(Z)=\big(l'(g(Z))\big)^{-1}.
\]
Differentiating the first identity at zero determines each next derivative of \(g\) from the lower derivatives, derivatives of \(l'\) and integer multiplicities, since \(l'(0)=1\).
The integral factor of the orthogonal chart therefore
retains integral derivatives in additive coordinates.
To compute its quadratic part on the selected matrices, use
the standardizing fourfold isometry of
Section~\ref{subsec:theta-normalization}. In the one-dimensional
coordinates the translated integral factor is supplied by
\eqref{eq:torsion-difference}. For \(Q\) modulo the period lattice the one-dimensional factor used in that calculation has the form
\[
  G_{jQ}(t)=e^{q_E(t)}(U_{jQ}\varsigma_*)(l(t)).
\]

The proof by reduction of the multiples \(jQ\) works also on the good curve at \(d\); the same gauge \(e^{q_E(t)}\) occurs for every \(j\). For a rational Hom direction
\(v\) with \(v\Gamma\) in the period lattice, the argument matrix
of the fourfold isometry has entries in the quadratic integers
localized at \(p\), by the multiplier-ring calculation of
Section~\ref{subsec:baseline-jets}. Thus the bilinear contraction
of two such argument columns is congruent modulo
\(\pi=\sqrt{-d}\) to their Hermitian contraction. Conjugating
one factor changes an integral quadratic-field entry by a
multiple of \(\pi\). Equivalently for two columns \(u,v\) in the orthogonal
coordinates,
\[
       u^{\mathsf T}v-u^*v=(u-\bar u)^{\mathsf T}v
                                 \equiv0\pmod{\pi}.
\]
The Hermitian contraction \(u^*v\) is the original trace
Hermitian value because the fourfold map is an isometry.
This computation also gives zero cross contractions modulo
\(\pi\) between disjoint copies. It remains to evaluate the
trace form on the directional tests in
\eqref{eq:unserved-normalized-frame}.

\proofstep{Separate the paired labels by integral selectors.}

Take a row of \(P_i^{-1}\) and a column of \(P_{i+1}\).
At label \(z\) multiply each by an integer of \(E\) that is
a unit there and clears the functional into the source dual
and the vector into the target. For a lattice \(M\), membership after localization gives exactly the clearing operation:
\[
  u\in M_{\mathfrak p}\quad\Longrightarrow\quad su\in M\text{ for some }s\in{\cal O}_E\setminus\mathfrak p.
\]
 This works on
any admissible pair agreeing with the selected pair locally.

If \(R,C\) denote the cleared row and column, \(a_l C R\) is a
global direction in the period module for every \(a_l\) from the
axis expression above. Evaluate that expression with coefficients
\(r_l/d_z\), then divide by the two clearing units in the
chosen embedding.
In fact the clearing factors \(s_1,s_2\) and the selected \(d_z\) are units, so
\[
  v_w\!\left(\frac{r_l}{s_1s_2d_z}\right)=v_w(r_l)\ge0 .
\]
 This is an integral combination giving exactly
one column of \eqref{eq:unserved-normalized-frame} and zero at
the other labels. All higher derivative remainders remain
integral after this linear substitution and products, by the
chain rule. Here and below \(d_z\) denotes the axis generator
at the type label of \(z\).

Expand the Hermitian trace on the two rank-one directions.
Each embedding contributes its row pairing times its column
pairing times \(a_l\overline{a}_{l'}\). The selector in the
first direction isolates the chosen embedding. For the second
direction the needed comparison is
\[
 \rho_{-,t}(\bar s)\equiv\rho_{+,t}(s)\pmod{\pi}
                  \qquad(s\in\mathcal O_E).
\]
Indeed the inertia involution at \(d\) has product labels
\(((1,-1),1)\). Composing the minus label with this involution
and complex conjugation gives the plus label at the same real
index \(t\). Inertia acts trivially on the residue field.
Its fixed part is unramified and the remaining quadratic field
is generated by \(\sqrt{-d}\). The extension over \(K\) is
unramified, so \(\pi\) remains a uniformizer.
Consequently the second selector kills all other partners
modulo \(\pi\). The source and target contractions on the surviving
pair give identity matrices modulo \(\pi\).

For instance if the
two row-clearing scalars are \(s_-,s_+\), the unspecialized
pairing is \(s_-\bar s_+\delta_{aa'}\), by the inverse-adjoint
identity for the two bases. After the two divisions the diagonal
entry is 
\[
  \rho_-(\bar s_+)/\rho_+(s_+)\equiv1\pmod\pi.
\]
The column contraction has the same calculation. Thus the
coefficient \(u_t\) in the surviving block is the scalar sum
of the products
\[
 \frac{r_{-,l}}{d_{-,t}}\,
 \rho_{-,t}(a_{-,l}\bar a_{+,l'})\,
 \frac{r_{+,l'}}{d_{+,t}},
\]
with the selectors for this pair of labels. Its terms are
integral because the inverse half different is integral here.
\proofstep{Extract the common quadratic phase from the fourth tensor.}

The opposite pairing gives the same coefficient modulo \(\pi\)
by symmetry of the original quadratic contraction. In particular
\(u_t\) can be fixed before the translates and their denominator
clearings. The quadratic sum on the four copies is therefore
the sum on each copy of
\[
        2\sum_{t,a,b}u_t X_{tab}Y_{tab}
\]
up to \(\pi\) times an integral quadratic polynomial. Because
\(v(S\pi)\ge0\) and two is a unit, the exponential of the error
multiplied by \(S/2\) has integral ordinary derivatives. This
follows also by the exponential product rule for derivatives
at zero.

Removing the common quadratic phase on the four
copies and taking the derivative at four equal indices bounds
the corresponding derivative on a single copy. Thus each
translate in \eqref{eq:unserved-reader} in its actual frame
is \(\exp(-S\sum u_tX_{tab}Y_{tab})\) times a series with
integral ordinary derivatives. The common exponential factors
out of the finite sum, including for the selected twist with
\(p\notin Q(\mu)\). This supplies the paired shape used in the
ramified bound. \label{cor:unserved-period-1}

\end{proof}

The base entry bound is zero at a level-free place away
from the discriminant whenever only integral jets are
used. This applies beyond the cutoff and at primes
with no applicable ordinary forced contribution. Section~\ref{subsec:baseline-jets} proves it in the actual
local bases, including at two and three. The stage bounds above are
uniform in the two classes of the entry, the word, and
each allowed tuple. Their inputs are the standard
theta charts of the same kernel datum, not charts
chosen separately for different coefficients of a jet.
\label{cor:unserved-period-2}
